\documentclass[12pt,a4paper]{article}

\usepackage[a4paper, margin=1in]{geometry}
\usepackage{longtable}
\usepackage{booktabs}
\usepackage{hyperref}
\usepackage{enumitem}
\usepackage{xcolor}

\hypersetup{
    colorlinks=true,
    linkcolor=blue,
    urlcolor=blue
}

\title{\textbf{Project Plan Document}\\
\vspace{0.3cm}
Hublet – ALPS\\
\vspace{0.2cm}
\large Architecting a Scalable Market Intelligence System}
\date{11 February 2026}

\begin{document}

\maketitle

\section*{Project Details}

\textbf{Project Number:} 46 \\[4pt]
\textbf{Client:} Vaibhav Singh, IIT Delhi \\[4pt]
\textbf{Mentor Email:} vaibhavsinghiitd28@gmail.com \\[6pt]

\textbf{Team Members:}
\begin{itemize}[noitemsep]
    \item Rachit Mehta – Backend System Developer
    \item Arijeet Paul – Backend \& Database Architecture
    \item Rakshit Agarwal – Data Engineering \& Database Management
    \item Jenik Gajera – Frontend \& QA Engineer
    \item Abhinav PV – DevOps Engineer
\end{itemize}

\section*{Problem Statement}

The Hublet project focuses on designing and implementing a scalable backend architecture for a real-estate market intelligence and lead-matching system.

Unlike traditional listing platforms, Hublet operates as a structured “black-box” matching engine that:

\begin{itemize}
    \item Ingests heterogeneous property data from web portals and APIs
    \item Parses unstructured buyer enquiries into structured intents
    \item Applies deterministic rule-based scoring for property matching
    \item Maintains a lead lifecycle state machine
\end{itemize}

The system emphasizes backend architecture, workflow orchestration, database design, and reliability. AI/ML-based valuation models are considered optional stretch components.

\section*{Technology Stack}

\subsection*{Backend}
\begin{itemize}
    \item Node.js with Express.js
    \item Prisma ORM
    \item SQLite (Development) → Supabase (Production)
    \item JWT Authentication
\end{itemize}

\subsection*{Frontend}
\begin{itemize}
    \item React 18+ with TypeScript
    \item Vite
    \item Axios
\end{itemize}

\subsection*{DevOps \& Tools}
\begin{itemize}
    \item Git \& GitHub
    \item GitHub Actions
    \item Docker \& Docker Compose
    \item Postman / Thunder Client
    \item Prisma Studio
\end{itemize}

\section*{Development Methodology}

The team follows a weekly milestone-driven development model with:

\begin{itemize}
    \item WhatsApp for daily coordination
    \item Google Meet for weekly review meetings
    \item GitHub for version control and CI/CD
    \item Markdown-based technical documentation
\end{itemize}

\newpage
\newpage
\section*{Milestone Schedule}

\renewcommand{\arraystretch}{1.4}

\subsection*{Release 1 (R1): Core Backend System}

\begin{center}
\begin{tabular}{|c|p{8cm}|c|c|}
\hline
\textbf{\#} & \textbf{Milestone} & \textbf{Date} & \textbf{Status} \\
\hline
1 & Team formation, client kickoff, and requirement gathering & 25/01/26 & Completed \\
\hline
2 & System architecture design and technology stack finalization & 29/01/26 & Completed \\
\hline
3 & Development environment setup and repository initialization & 30/01/26 & Completed \\
\hline
4 & Database schema design and Prisma configuration & 01/02/26 & Completed \\
\hline
5 & CRUD API implementation (Buyer, Seller, Property) & 02/02/26 & Completed \\
\hline
6 & Keyword-based intent parsing module implementation & 04/02/26 & Completed \\
\hline
7 & Web scraping pipeline setup and cron scheduler integration & 10/02/26 & Completed \\
\hline
8 & Admin, Buyer, and Seller dashboard (basic views) & 08/02/26 & Completed \\
\hline
9 & Rule-based weighted matching engine implementation & 12/02/26 & In Progress \\
\hline
10 & Lead lifecycle state machine implementation & 12/02/26 & In Progress \\
\hline
11 & JWT authentication and protected route middleware & 13/02/26 & Pending \\
\hline
12 & \textbf{R1 Deliverable: Requirements Document Submission} & 14/02/26 & Completed \\
\hline
\end{tabular}
\end{center}

\vspace{1cm}

\subsection*{Release 2 (R2): Scalability, Security \& Deployment}

\begin{center}
\begin{tabular}{|c|p{8cm}|c|c|}
\hline
\textbf{\#} & \textbf{Milestone} & \textbf{Date} & \textbf{Status} \\
\hline
13 & Input validation (Zod) and scraped data normalization & 20/02/26 & Pending \\
\hline
14 & Role-Based Access Control and security middleware integration & 22/02/26 & Pending \\
\hline
15 & API documentation using Swagger / OpenAPI & 23/02/26 & Pending \\
\hline
16 & Unit and integration testing (60\%+ coverage) & 25/02/26 & Pending \\
\hline
17 & Docker setup and containerized local environment & 27/02/26 & Pending \\
\hline
18 & Asynchronous job queue integration (BullMQ) & 03/03/26 & Pending \\
\hline
19 & Email notification system and workflow alerts & 05/03/26 & Pending \\
\hline
20 & Advanced filtering, deduplication, and analytics dashboard & 08/03/26 & Pending \\
\hline
21 & Monitoring, logging, and performance optimization & 12/03/26 & Pending \\
\hline
22 & CI/CD setup and production deployment configuration & 18/03/26 & Pending \\
\hline
23 & Final QA, security audit, and system integration testing & 22/03/26 & Pending \\
\hline
24 & \textbf{R2 Deliverable: Final Demo and Deployment Documentation} & 25/03/26 & Pending \\
\hline
\end{tabular}
\end{center}


\section*{Stretch Goals (Post R2)}

\begin{itemize}
    \item ML-based property valuation (Scikit-learn)
    \item Collaborative filtering for matching
    \item A/B testing: ML vs Rule-based engine
    \item Deploy ML microservice (Flask/FastAPI)
\end{itemize}

\section*{Risk Assessment}

\begin{itemize}
    \item Web scraping reliability issues
    \item Data inconsistency across portals
    \item Authentication security vulnerabilities
    \item Performance degradation with scale
\end{itemize}

Mitigation strategies include input validation, indexing, logging, CI/CD testing, and modular service separation.

\section*{Conclusion}

This project plan outlines the structured development of Hublet as a scalable backend system for real estate intelligence and deterministic lead matching. The focus remains on backend engineering, workflow orchestration, and system reliability, with AI-based components treated as optional extensions.

\vfill
\begin{center}
\textbf{Submitted on: 12 February 2026}
\end{center}

\end{document}
